\documentclass[11pt,letterpaper]{article}
\usepackage[margin=0.8in]{geometry}
\usepackage[T1]{fontenc}
\usepackage{lmodern}
\usepackage{amsmath,amssymb,xcolor,listings,enumitem,hyperref}
\hypersetup{colorlinks=true,linkcolor=blue,urlcolor=blue,plainpages=false,pdftitle={CMSC433 Topic Practice}}
\definecolor{paper}{RGB}{246,248,250}
\lstset{basicstyle=\ttfamily\small,backgroundcolor=\color{paper},columns=fullflexible,
keepspaces=true,breaklines=true,showstringspaces=false,frame=single,
rulecolor=\color{gray!30},aboveskip=0.5em,belowskip=0.5em}
\lstnewenvironment{codeanswer}{\lstset{backgroundcolor=\color{white},frame=single}}{}
\setlist[enumerate]{itemsep=0.7em}
\setlength{\parindent}{0pt}
\setlength{\parskip}{0.5em}
\ExplSyntaxOn
\NewDocumentCommand{\code}{m}{\texttt{%
\tl_set:Nx\l_tmpa_tl{\tl_to_str:n{#1}}%
\seq_set_split:NnV\l_tmpa_seq{~}\l_tmpa_tl
\seq_map_indexed_inline:Nn\l_tmpa_seq{%
\int_compare:nNnT{##1}>{1}{\penalty0\relax\space}%
\str_map_inline:nn{##2}{####1\kern0pt\relax}}}}
\ExplSyntaxOff
% TYPE YOUR WRITTEN ANSWERS INSIDE THESE BRACES in the included lesson.tex files.
\newcommand{\answer}[1]{\par\nopagebreak[4]\smallskip\noindent\textbf{Answer: }#1\par\medskip}
\newcommand{\task}[1]{\subsection*{#1}}
\newcommand{\source}[1]{{\small\textbf{Source:} #1}\par}
\newcommand{\ready}[1]{{\small\textbf{Checkpoint:} #1}\par}

\raggedbottom
\newcommand{\sheettitle}[3]{\noindent{\small CMSC433 / Part #1}\par
{\Large\bfseries #2}\par{\large #3}\par\smallskip\hrule\medskip}
\newcommand{\mockrules}{{\small Topic mock, using the sample's question styles.
Attempt independently; feedback after the whole file. Untimed practice, no points
assigned. Current exam scope and rules are unconfirmed.}\par}

\begin{document}
\sheettitle{02}{Functor, Applicative, and State types}{Guided practice}

\textbf{Use with the chat:} work through one question at a time.

\textbf{Idea.} \code{fmap} changes the value inside a context. For
\code{State s a}, \code{s} is the stored type and \code{a} is the returned type.

\textbf{Worked example.} \code{null <$> get} has type \code{State [a] Bool}:
\code{null :: [a] -> Bool} forces the stored state to be a list; the result
is a Boolean and the State context stays in place.

\task{1. A function inside Maybe}
Infer \code{(,) <$> Just False}. First write the type of \code{(,)} after
supplying a Boolean, then retain the outer context.
\answer{Intermediate function :: [ ]; whole expression :: [ ].}

\task{2. Apply a contextual function}
Infer \code{(++) <$> Just "a" <*> Just "bc"}.
Give the type immediately before \code{<*>} and the whole type.
\answer{[Intermediate type, final type, and why <*> fits.]}

\task{3. Track the state separately}
Infer \code{get >>= \xs -> put (reverse xs)}.
\answer{Stored type: [ ]; returned type: [ ]; whole type: [ ].}

\task{4. Changed independent question}
Infer \code{\f -> modify f >> get}. Explain the relationship between the
input and output types of \code{f}.
\answer{[Most general type and reasoning.]}
\source{Sample Q1, p. 1; Monads Functor/Applicative section; State interface.}

\end{document}
